% Part: applied-modal-logics
% Chapter: epistemic-logic
% Section: relational-models

\documentclass[../../../include/open-logic-section]{subfiles}

\begin{document}

\olfileid{aml}{el}{rel}

\olsection{تعلقی نمونے}

علم دیاں منطقاں دا بنیادی معنیاتی خیال عام امکانی منطقاں والا اے۔
تعلقی نمونیاں وچ ہُن وی جہاناں دا سیٹ تے اک قدر بندی اے جیہڑی طے کردی
اے کہ کہڑے !!{propositional
variable}s کہڑے جہاناں اُتے «سچ» نیں۔ جے اسیں صرف اک فاعل دی گل
کر رہے آں، تاں معمول وانگ رسائی دا اکّو تعلق اے۔ کئی فاعلاں والی علم
دی منطق وچ رسائی دے تعلقاں دا سیٹ ہوندا اے، فاعلاں دے نشاناں دے
سیٹ وچ ہر~$a$ لئی اک۔ اوس سیٹ دا ناں~$G$ اے۔

اک \emph{تعلقی نمونے} وچ جہاناں دا سیٹ اے، جیہڑے دو رکنی رسائی
تعلقاں نال جڑے نیں، ہر فاعل لئی اک تعلق۔ نال اک قدر بندی اے جیہڑی
طے کردی اے کہ کہڑے !!{propositional variable}s کہڑے جہاناں اُتے سچ نیں۔

\begin{defn}
  کئی فاعلاں والی علم دی زبان دا اک \emph{نمونہ} ترتیبی تِکڑی
  $\mModel{M} = \tuple{W, R, V}$ اے، جتھے
  \begin{enumerate}
  \item $W$ «جہاناں» دا اک غیر خالی سیٹ اے۔
  \item ہر $a \in G$ لئی، $ {R}_a$،~$W$ اُتے رسائی دا دو رکنی تعلق اے۔
  \item $V$ اوہ فنکشن اے جیہڑا ہر !!{propositional
    variable}~$p$ نوں ممکن جہاناں دا سیٹ $V(p)$ دیندا اے۔
  \end{enumerate}
  جدوں $R_a ww'$ ہووے، اسیں آکھدے آں کہ $w'$، \emph{فاعل a لئی
  رسائی وچ اے}~$w$ توں۔ جدوں $w \in V(p)$ ہووے، اسیں آکھدے آں کہ
  $p$،~$w$ اُتے \emph{سچ اے}۔
\end{defn}

کم کرن دا طریقہ معیاری امکانی منطق والا اے، بس رسائی دے ہور تعلق
شامل نیں۔ کسے فاعل لئی اسیں عام طور اُتے اوہدے رسائی دے تعلق نوں
اوہدی جانکاری دی حالت بارے کجھ دسن والا سمجھاں گے۔ مثلاً اسیں اکثر
$R_a ww'$ دا مطلب ایہ لیندے آں کہ $w'$،~$a$ دی جانکاری نال میل
کھاندا اے جیہڑی اوہنوں~$w$ اُتے اے۔ دوجے لفظاں وچ،~$w$ اُتے اوہ
جہان $w$ تے جہان~$w'$ وچ فرق نہیں دس سکدا۔

\end{document}
